\printedpage{14}
% head: A. Grothendieck : Dérivateurs, Ch. IX, Édition des Universités de Montpellier II et Paris VII
% foot: Version 12 mars 2002

The existence of
\[
\begin{aligned}
\Omega :\ & \mathcal{M}W^{-1} \ \longrightarrow\ \mathcal{M}W^{-1}, \quad \text{or of}\\
\Sigma :\ & \mathcal{M}W^{-1} \ \longrightarrow\ \mathcal{M}W^{-1},
\end{aligned}
\]
or even of both at once, appears as an extremely strong hypothesis \cercle{a} on $(\mathcal{M}, W)$. This condition \textit{plus} the one \cercle{b} that the $\Chu_{\infty}(a,b)$ be homotopic to Thomason-Kan categories, or (which comes to the same thing, more or less?) that $Y \longmapsto \mathrm{Hom}_{\mathcal{M}(Y)W_Y^{-1}}(a_Y, b_Y)$ transform Hot-equivalences of Cat into bijections (even if only for fixed $Y$ \dots), seems the crucial test for knowing whether ``$(\mathcal{M}, \Sigma)$ \texttt{[rather} $(\mathcal{M}, W)$\texttt{]} lends itself to a homotopy theory ``in $\mathcal{M}$'' ''. I doubt that

\msp{page 23}

these conditions can be satisfied for finite $\mathcal{M}$, except in the trivial case where $\mathcal{M} = \emptyset$, where $\mathcal{M}\Sigma^{-1}$ \texttt{[rather} $\mathcal{M}W^{-1}$\texttt{]} is equivalent to the one-point category.

A third test \cercle{c} which seems important to me (but perhaps this follows from the rest), is the formulas in $\Hot_\bullet$
\[
\left\{
\begin{aligned}
\Homext_\bullet(\Sigma a, b) \ &\simeq\ \Homext_\bullet(a, \Omega b)\\
&\simeq\ \Omega\Homext_\bullet(a, b),
\end{aligned}
\right.
\tag*{$(*)$}
\]
where $\Homext_\bullet$ denotes the external Hom with values in $\Hot_\bullet$. One must view it separately (assuming the existence either of $\Omega$, or of $\Sigma$) as two formulas ``dual'' to one another,
\[
\left\{
\begin{aligned}
\Homext_\bullet(a, \Omega b) \ &\simeq\ \Omega\Homext_\bullet(a, b)\\
\Homext_\bullet(\Sigma a, b) \ &\simeq\ \Omega\Homext_\bullet(a, b),
\end{aligned}
\right.
\tag*{$(**)$}
\]
isomorphisms of \textit{groups} in $\Hot_\bullet$. (\textbf{NB} $\Omega b$ is a group object of $\mathcal{M}W^{-1}$, for it represents a contrafunctor with values in Gr, and likewise $\Sigma a$ is a cogroup of $\mathcal{M}W^{-1}$. The functor $y \longmapsto \Homext(a, y)$ must commute with finite products, and thus transform group objects into likewise. And likewise $x \longmapsto \Homext(x, b)$ must transform finite sums into finite products, hence cogroups into groups (\footnote[8]{The functor $\Omega$ must come from an endomorphism of the derivator $\mathbf{D} = \mathbf{D}_{\mathcal{M},W}$, and the latter must be left exact, and dually for $\Sigma$.}).)

I have convinced myself that the isomorphisms $(**)$ have a chance of existing only if one has the prerequisite condition \cercle{b}.

\newpage

\printedpage{15}
% head: A. Grothendieck : Dérivateurs, Ch. IX, Édition des Universités de Montpellier II et Paris VII
% foot: Version 12 mars 2002

\msp{page 24}

\section*{3\quad Attempt at constructing the functor $\Omega$}

There is probably little chance of proving anything about the functor
\[
\Omega :\ \mathcal{M}W^{-1} \longrightarrow \mathcal{M}W^{-1},
\]
without having in hand a more or less explicit \textit{construction} of it. I am going to try, taking my inspiration from \textsc{Quillen}, to work out general conditions on $(\mathcal{M}, W)$ (more general than those of \textsc{Quillen}) which ensure the existence of $\Omega$, and allow one to construct it. The geometric idea goes back to \textsc{Serre-Cartan}. The task is, for an object $S$ of $\mathcal{M}$, to construct an object $\mathrm{Ch}(S)$ \textit{in $\mathcal{M}$,} playing the role of a ``path space'', and which must fit into a factorization diagram for $\mathrm{diag}_S : S \longrightarrow S \times S$ (we assume that in $\mathcal{M}$ finite $\varprojlim$, hence finite products, exist)
\[
\begin{tikzcd}
S \arrow[d, "i_S"] \arrow[dd, bend right=40, "\mathrm{diag}_S"'] & \\
\mathrm{Ch}(S) \arrow[d, "p_S"] & \\
S \times S & e. \arrow[l]
\end{tikzcd}
\]
The essential property is that $p$ be a \textit{``fibration''} (in a suitable sense, to be worked out).

\msp{page 25}

As a subsidiary matter it may be useful (perhaps) that $i$ be a ``cofibration'' (\footnote[9]{and perhaps that $S \longrightarrow e$ also be a fibration?}). One defines $\Omega(S)$ as the fiber of $\mathrm{Ch}(S)$ over $S{\times}S$ (with respect to the unique morphism $e = e_{\mathcal{M}} \longrightarrow S\times S$, where $e_{\mathcal{M}}$ is the zero object of $\mathcal{M}$).
\[
\begin{tikzcd}[column sep=large]
\mathrm{Ch}(S) \arrow[d] & \Omega S \arrow[l] \arrow[d]\\
S \times S & e. \arrow[l]
\end{tikzcd}
\]

\newpage

\printedpage{16}
% head: A. Grothendieck : Dérivateurs, Ch. IX, Édition des Universités de Montpellier II et Paris VII
% foot: Version 12 mars 2002

This is an object of $\mathcal{M}$, which we shall regard mainly as an object of $\mathcal{M}W^{-1}$. We must make sure that we have a bijection functorial in $T$ in $\mathcal{M}W^{-1}$
\[
\underbrace{\pi_1\Chu^W_\infty(T, S)}_{\overset{\text{def}}{=}\,\pi_1(T,S)} \ \simeq\ \underbrace{\mathrm{Hom}_{\mathcal{M}W^{-1}}(T, \Omega S)}_{\overset{\text{def}}{=}\,\pi_0\Chu^W_\infty(T,\Omega S)}.
\]
For this, one must compare
\[
\Chu^W_\infty(T, \Omega S) \qquad\quad \text{and} \qquad\quad \underbrace{\Omega_{\Hot_\bullet}(\Chu^W_\infty(T, S))}_{\overset{\substack{\text{``by definition''}\\ \text{for $\Omega$ of $\mathrm{Hot}_\bullet$}}}{=}\ \Chu_\infty(\mathrm{Ch}^W_\infty(T,S);0,0)} \qquad ,
\]
since it is a matter of constructing a bijection between their $\pi_0$.